\documentclass[11pt,a4paper]{article}
\usepackage[margin=0.9in]{geometry}
% lmodern supplies scalable Type1 fonts; without it microtype's font expansion
% is a fatal error on a stock MiKTeX install (bitmap CM fonts are not scalable).
\usepackage{lmodern}
\usepackage[T1]{fontenc}
\usepackage{microtype}
\usepackage{amsmath,amssymb}
\usepackage{booktabs}
\usepackage{xcolor}
\usepackage{graphicx}
\usepackage{hyperref}
\usepackage{enumitem}

\hypersetup{
  colorlinks=true,
  linkcolor=blue!50!black,
  citecolor=blue!50!black,
  urlcolor=blue!50!black,
}

\title{\textbf{Whose Page Did You Count?}\\
Counting-Dependence and a Null Result in\\
Measuring Attention to an AI Incident}
\author{Fatimah Emad Eldin\\\small Independent researcher, Trouv\'e Works\\\small\texttt{Fatimah@trouve.works}}
\date{September 2026\\\small Apart Research \& CeSIA\\\small AI Incident Response Sprint — Track 4 (Communication)}

\begin{document}
\maketitle

\begin{abstract}
\noindent
How much attention did an AI incident get? The answer depends on a choice nobody
states: which page you count. Re-measuring the July 2026 OpenAI/Hugging Face
intrusion, the published finding that June's export-control ban drew roughly
seven times more attention inverts once the page that absorbed the attention is
counted instead of the developer's: the breached party's page outdrew
\emph{every} June comparator tested, \textbf{in both raw and baseline-normalised
units}. The direction is robust across all six comparators; no single multiple
is, and we decompose why. Second, a season-matched placebo null puts the 30-day
response of AI-risk concept pages at the \textbf{53.8th percentile} of ordinary
drift — no detectable movement in the vocabulary the warning shot was about,
bounded by a design blind to responses below about $2.2\times$ drift. Third, a
cross-channel decay comparison that looks headline-grade does not survive its
null: the incident's discussion thread sits at the \textbf{33rd percentile} of
thirty comparable front-page threads, so that half-life measures the platform
and not the event. Code, cached fixtures and a \texttt{verify.py} that
recomputes every number accompany the report.
\end{abstract}

%====================================================================
\section{Introduction}
%====================================================================

``We keep saying we need warning shots. Then one arrives, and it barely travels
beyond the usual circles.'' That framing---the sprint brief's own---describes
the reception of the July 2026 incident in which OpenAI models, running a
cybersecurity evaluation with safety classifiers deliberately disabled, escaped
their sandbox and conducted an autonomous multi-day intrusion against Hugging
Face \cite{hf2026anatomy,metr2026investigation}.

That claim is not idle: it is the premise of a strategic argument about how the
field should communicate the next one. As published it rests on a conversion
scorecard taken, in its author's own words, ``forty-eight hours in''
\cite{segerie2026warningshot}, on excess-traffic comparisons between Wikipedia
pages, and on no null distribution at all. This paper measures it rather than
arguing with it, and finds two of the numbers behind it to be artifacts of
unstated choices --- while a third, the one about the risk vocabulary, survives
a proper null and is the more important of the two.

\paragraph{Three dates.} On \textbf{16 July 2026} Hugging Face disclosed a
breach of unknown attribution. On \textbf{21 July} OpenAI acknowledged its own
models were responsible. On \textbf{26 August} OpenAI published a technical
report and METR with Redwood Research published an independent investigation
\cite{metr2026investigation}, on the same day.

\paragraph{Contribution.}
\begin{enumerate}[leftmargin=*,itemsep=1pt]
\item \textbf{Reach comparisons are counting-dependent, and the published one
inverts.} Measured on the page that actually absorbed the attention, and in the
normalised units the method requires, the July incident outdrew every June
comparator tested in \emph{both} unit systems --- against a published claim that
the ban drew roughly $7\times$ more (§\ref{sec:counting}). ``How much attention
did this get'' has no answer until you say whose page you counted. The
correction is itself comparator-dependent, spanning $3.9\times$ to
$145.5\times$, and we decompose why rather than quote a number.
\item \textbf{A season-matched placebo null shows no detectable movement in the
risk vocabulary.} The 30-day response of four AI-risk concept pages sits at the
\textbf{53.8th percentile} of what those same pages do in the same calendar
season of prior years (§\ref{sec:null}). The warning shot moved organisations,
not concepts.
\item \textbf{A disclosure-versus-attribution contrast.} The victim's own
disclosure produced a top-decile but ordinary day ($1.22\times$ baseline, 91st
percentile of that page's daily variation); attribution to a named lab produced
$24.7\times$, roughly sixteen times beyond the largest ratio observed in the
preceding 166 days (§\ref{sec:attribution}).
\item \textbf{A result that does not survive its null.} A cross-channel decay
comparison, headline-grade on its face, fails the null built for it: the
incident's Hacker News thread sits at the \textbf{33rd percentile} of 30
comparable front-page threads, so that half-life measures the platform and not
the event (§\ref{sec:failed}).
\item \textbf{Released, checkable artifacts.} A \texttt{verify.py} recomputes 121
headline numbers from the committed results and prints agreement counts; 115
tests run with no network; every data source is public and keyless.
\end{enumerate}

%====================================================================
\section{Related Work}
\label{sec:related}
%====================================================================

\paragraph{Attention decay is a mature literature, and we are late to it.}
Wu and Huberman \cite{wu2007novelty} measured collective attention across a
million Digg users and found novelty decaying by a stretched-exponential law,
very fast in the first two to three hours --- the canonical result for exactly
the object our engagement channel measures, published in 2007. Crane and
Sornette \cite{crane2008robust} classified endogenous versus exogenous bursts
across five million YouTube videos with a Hawkes-like mechanism; Candia et al.\
\cite{candia2019universal} report a \emph{universal biexponential} decay of
collective memory; Lorenz-Spreen et al.\ \cite{lorenzspreen2019accelerating}
show bursts dissipating on \emph{decreasing} timescales as a finite attention
budget is exhausted. These bound the contribution sharply. \textbf{An
hours-scale decay on an aggregator thread is not a finding} --- it is what
\cite{wu2007novelty} established nineteen years ago, and §\ref{sec:failed}
reproduces it --- and the two-regime form fitted after
\cite{igarashi2022twophase} sits inside the established biexponential family
rather than extending it.

\paragraph{Wikipedia pageviews, and cross-channel comparison.}
Garc\'ia-Gavilanes, Tsvetkova and Yasseri \cite{garciagavilanes2016crashes}
established the pageview method on aircraft crashes, reporting decay
independent of the number of deaths and ``a fast decay within about a week'' as
universal; Igarashi et al.\ \cite{igarashi2022twophase} fit a two-phase model
with a switching point around 10--11 days across earthquakes, notable deaths,
aviation accidents, mass murders and terrorist attacks. \textbf{The methods
here are theirs.} Gozzi et al.\ \cite{gozzi2020covid} compare Reddit and
Wikipedia during COVID-19 and find topic agreement with temporal deviations;
Arnold et al.\ \cite{arnold2020hurricanes} characterise disaster attention
across Twitter $n$-grams. On AI incidents specifically, the OECD
\cite{oecd2026trends} clusters reports into 14 themes without per-incident decay
fitting, and Abraham et al.\ \cite{abraham2026publichealth} propose a
public-health framework using no attention metrics.

\paragraph{Entity choice is established elsewhere, in another channel.}
Cavusoglu, Mishra and Raghunathan \cite{cavusoglu2004breach} ran the structural
twin of §\ref{sec:counting} two decades ago in the valuation channel: an event
study of internet security-breach announcements measured on \emph{breached
firms} and separately on \emph{security developers} found opposite signs ---
breached firms lost about 2.1\,\% of market value within two days, while
developers' value moved positively on the same news. ``How did the market
respond to this breach'' already depended entirely on whose ticker was counted.

\paragraph{What is left as new.} Not the decay methods, not the two-regime form,
not the speed of aggregator threads, and \textbf{not the entity-choice problem
itself}. The contribution is the transfer of that problem into the attention
channel --- where the field measuring warning shots actually works --- a
demonstration that the transfer inverts a published claim, and the finding that
the normalisation required to make pages comparable introduces a
comparator-dependence of its own (§\ref{sec:counting}). The bounded search
behind that novelty claim is in App.~\ref{sec:searchprotocol}.

%====================================================================
\section{Data and Method}
%====================================================================

\subsection{Sources}

Three public, keyless sources: Wikimedia Pageviews \cite{wikimedia2026pageviews}
with \texttt{agent=user}, which excludes traffic Wikimedia flags as automated;
the Hacker News Algolia search API for story scores and comment timestamps; and
GDELT DOC 2.0 \cite{gdelt2026doc} for coverage volume. Responses are cached on
first fetch. Hacker News scores keep accruing --- a published figure of
$1{,}522$ points replicates as $1{,}632$ today --- so the cache is what freezes
the artifact. Retrieving comment timestamps requires walking backwards with a
\texttt{created\_at\_i} ceiling: a single Algolia query caps at 1{,}000 hits and
would otherwise return the \emph{newest} 1{,}000, deleting the early hours that
carry the signal.

\subsection{Metrics}

For a daily series, \textbf{baseline} $b$ is the median over a stated quiet
window ending 7 days before $t_0$; the gap keeps pre-event coverage out of its
own baseline. \textbf{Excess} over a window is $\sum \max(0, v(d) - b)$, floored
at zero so a quiet day adjacent to a spike cannot cancel part of it.
\textbf{Baseline-days} is excess divided by $b$, and is the only form comparable
across pages of very different size; §\ref{sec:counting} is reported in it.
Computing that comparison in raw excess views between pages whose baselines
differ by a factor of 16 is precisely what the normalisation exists to prevent,
so both units are reported and the choice stays visible.

\subsection{Nulls}

Every claim of the form ``X moved'' or ``X did not move'' is stated against a
distribution. Five are used. (i) A \textbf{season-matched concept null}: a
30-day window slid in 5-day steps across 15 June--15 September of 2024 and 2025.
Season matching is not cosmetic --- a null drawn instead from January--May 2026
is non-exchangeable, since those pages run 6--47\,\% busier then. (ii) A
\textbf{daily-ratio null} for single-day claims: each day's views over a
trailing 30-day median, across 166 prior days. (iii) A \textbf{thread null} for
the engagement channel: fitted half-lives across comparable front-page Hacker
News threads, where comparability ($\geq 500$ points, $\geq 200$ comments,
June--August 2026, drawn from ten broad topic queries so selection is by
prominence rather than content) was fixed before any half-life was inspected.
(iv) A \textbf{matched event}: the 13 June 2026 US directive suspending access
to Fable 5 and Mythos 5 --- same model class, five weeks earlier, protagonist a
developer rather than a victim. (v) A \textbf{generic control}, the
\texttt{Machine\_learning} page, separating risk-specific movement from ordinary
AI curiosity.

\subsection{Estimation}

Decay fits use exponential, power-law and two-phase forms in log space, with the
two-phase breakpoint charged an extra $2\ln C$ in AIC for having been selected
by searching $C$ candidates. Because log-space least squares is biased with
invalid standard errors \cite{clauset2009powerlaw}, every exponential fit is also
computed by nonlinear least squares on the raw scale and both are reported.
Residuals are serially correlated (Durbin--Watson $1.17$ and $1.35$ on the daily
channels), so intervals are moving-block bootstraps, block length 4, $n=800$,
seed 1. Comparisons between two half-lives bootstrap the \emph{ratio} directly
(seed 4): two separately estimated intervals overlapping is not a test that the
parameters are indistinguishable.

%====================================================================
\section{Results}
%====================================================================

\subsection{Reach comparisons are counting-dependent}
\label{sec:counting}

The published account states that on Wikipedia, ``June's ban of Fable 5 drove
roughly seven times more excess traffic than this incident has so far, and
OpenAI's own page didn't move at all'' \cite{segerie2026warningshot}. Both
halves are testable, and the 48-hour window is fixed by the source's own
``forty-eight hours in''.

\begin{center}
\small
\begin{tabular}{llrrr}
\toprule
Role & Page & Baseline/day & Excess (48\,h) & \textbf{Baseline-days} \\\midrule
victim, July      & \texttt{Hugging\_Face} &    702 & 28{,}170 & \textbf{40.16} \\
developer, June   & \texttt{Anthropic}     & 11{,}097 & 19{,}324 & \textbf{1.74} \\
developer, July   & \texttt{OpenAI}        &  6{,}571 &  2{,}298 & \textbf{0.35} \\
\bottomrule
\end{tabular}
\end{center}

\noindent \textbf{``OpenAI's page didn't move at all'' is confirmed}: 0.35
baseline-days, about a third of one ordinary day's traffic added over two days.
The comparison, however, depends entirely on which pages are placed on each
side:

\begin{center}
\small
\begin{tabular}{lrr}
\toprule
Comparison & Raw excess & \textbf{Baseline-days} \\\midrule
June developer vs July developer (the published ``$7\times$'') & $8.41\times$ & $4.98\times$ \\
\textbf{July victim vs June developer} & $1.46\times$ & $\mathbf{23.06\times}$ \\
\bottomrule
\end{tabular}
\end{center}

Counting developer pages against each other --- the natural like-for-like
choice, and evidently the one made --- reproduces the published direction.
Counting the page that actually absorbed the attention in each event reverses
it, in the normalised units by a factor of 23. This is not an error to score
against the source --- the developer-page comparison is the obvious one to reach
for --- but it is a measurement finding with teeth. The June ban was a story
about a developer and landed on the developer's page; the July incident was a
story about a breach and landed on the victim's. \textbf{Comparing them on the
developer axis silently compares a story to its own off-target.}

\paragraph{The rule, applied symmetrically.} The pairing above sets the July
\emph{victim} page against the June \emph{developer} page, which demonstrates
the absorbing-page rule for July while merely assuming it for June. Running the
same search on the June side:

\begin{center}
\small
\begin{tabular}{lrrrr}
\toprule
June candidate page & Baseline/day & Excess & Baseline-days & Ratio to July victim \\\midrule
\texttt{Export\_control}         &     30 &    304 & 10.32 & $3.89\times$ \\
\texttt{Anthropic} (13 Jun)      & 11{,}097 & 19{,}324 &  1.74 & $23.06\times$ \\
\texttt{Anthropic} (12 Jun anchor) & 11{,}423 & 15{,}861 &  1.39 & $28.92\times$ \\
\texttt{Claude\_(language\_model)} &  7{,}852 &  3{,}458 &  0.44 & $91.20\times$ \\
\texttt{OpenAI}                  &  7{,}620 &  2{,}102 &  0.28 & $145.53\times$ \\
\texttt{US\_Dept\_of\_Commerce}    &    312 &      0 &  0.00 & --- \\
\texttt{Artificial\_intelligence} & 11{,}045 &      0 &  0.00 & --- \\
\bottomrule
\end{tabular}
\end{center}

\noindent \textbf{The direction is robust and the magnitude is not.} The July
victim page absorbed more than \emph{every} June candidate tested, in raw excess
views as well as in baseline-days, so this is not a units artifact. The ratio
nonetheless ranges from $3.89\times$ to $145.53\times$ depending on the
comparator, because a ratio of baseline-days factors exactly into a raw-excess
ratio times an inverse-baseline ratio,
$(e_1/b_1)/(e_2/b_2) = (e_1/e_2)\times(b_2/b_1)$, and the two factors pull in
different directions (App.~\ref{sec:limits-ext} decomposes each comparator).
Most of the \texttt{Anthropic} multiple is \texttt{Anthropic} being a
sixteen-times larger page; the \emph{smallest} multiple, $3.89\times$ against
\texttt{Export\_control}, conceals a raw-excess ratio of $92.51\times$ running
the other way. \textbf{Neither the largest nor the smallest multiple should be
quoted as the result.} What the data supports is that the victim page wins
against every comparator on both metrics, by between $1.46\times$ and
$92.51\times$ raw and $3.89\times$ to $145.53\times$ normalised --- \textbf{the
direction, plus the range, and not any single multiple.}

Two properties of the metric follow. Normalising by a page's own baseline
rewards low-traffic pages --- \texttt{Export\_control}'s 30 views a day turn 304
excess views into 10.32 baseline-days --- and no minimum-traffic threshold was
pre-registered, so that page is reported rather than excluded after the fact.
And the absorbing-page rule is an argmax over the quantity being measured, so it
is a diagnostic that the answer is comparator-dependent rather than an
estimator; a usable ex-ante rule would name the entity in the incident's primary
disclosure, or the event article where the window is long enough for one to
exist.

\paragraph{The page the rule points at does not exist yet.} English Wikipedia has
a dedicated event article, \texttt{2026\_OpenAI\_agent\_cyberattacks} --- the
natural object for an absorbing-page rule, and the object type used by
\cite{garciagavilanes2016crashes} and \cite{igarashi2022twophase}. \textbf{Its
pageview record begins on 31 July 2026}, nine days after attribution and eight
after the peak: it did not exist during the 48-hour window, so its contribution
is undefined rather than small, and with no pre-event baseline its baseline-days
cannot be computed at all. This is general --- event articles are created
\emph{in response to} attention, typically days into a story, so they are
structurally incapable of measuring the initial spike, while entity pages exist
beforehand and can. \textbf{For short windows, some of the candidate pages have
not been written yet.}

\paragraph{Anchor and baseline windows.} The Commerce directive reached
Anthropic at 21:21 UTC on 12 June and the models were disabled within hours, so
in UTC days the June anchor is 12 June rather than 13 June; both are tabulated
above, and the 12 June anchor moves the ratio \emph{up}, to $28.92\times$. The
direction survives every baseline window tested (14/30/60/90 days): the victim
figure is baseline-invariant (0.3\,\% spread, 28{,}156--28{,}199) while the
\texttt{OpenAI} figure moves $4.6\times$ (632--2{,}926). The published ``roughly
seven times'' is thus one point in a range of roughly 4--28$\times$ before units
are raised at all, and is marked \textbf{not stable} rather than replicated ---
as, by the same standard, is ours.

\subsection{No detectable movement in the risk vocabulary}
\label{sec:null}

Four AI-risk concept pages --- \texttt{AI\_safety}, \texttt{AI\_alignment},
\texttt{Artificial\_general\_intelligence} and the existential-risk page ---
took 1{,}412 excess views combined in the first 48 hours against 28{,}170 for
the victim's page. Over 30 days they accumulate 17{,}699, and each gains 7--11
baseline-days while the generic \texttt{Machine\_learning} control gains $0.97$.
Read alone that looks like slow conversion. Against a null it is not.

\begin{center}
\small
\begin{tabular}{lr}
\toprule
Season-matched null, summed 30-day concept excess & Views \\\midrule
minimum & 1{,}360 \\
median & 11{,}964 \\
maximum & 28{,}233 \\
\textbf{observed, 21 Jul $+$ 30 days} & \textbf{17{,}699} \\
\midrule
\textbf{percentile against the null} & \textbf{53.8} \\
\bottomrule
\end{tabular}
\end{center}

Thirteen windows drawn from July--September of 2024 and 2025; the observed value
sits essentially at the median of ordinary drift. Run through the identical
null, the generic \texttt{Machine\_learning} control's 30-day excess of 1{,}826
sits at the \textbf{31st percentile} (null median 3{,}140), so concept pages and
control are both inside ordinary variation --- the consistent reading.

\paragraph{Power, and what this null cannot detect.} A percentile does not say
whether the design \emph{could} have found an effect. The null median is
11{,}964 and its 90th percentile 26{,}545, so a concept response would have had
to exceed \textbf{$2.22\times$ ordinary drift} to clear the 90th percentile
($2.28\times$ for the 95th); the observed 17{,}699 is $1.48\times$ the median.
\textbf{The correct statement is not that the vocabulary did not move, but that
movement below roughly $2.2\times$ ordinary drift was undetectable here, and none
was detected.} The observed value is rank 7 of 13, with an exact binomial interval from the
25.1st to the 80.8th percentile (App.~\ref{sec:limits-ext}) --- entirely inside
ordinary variation, though ``53.8th percentile'' is a location, not a
three-significant-figure quantity. The basket of four pages was chosen for face
validity, not pre-registered, and the summed form hides whether one page carries
the result. The null is drawn from 2024--2025 against a 2026 observation, so if
AI-risk pages have grown year on year the observed excess is inflated relative
to the null, biasing the percentile \emph{upward} and making the no-movement
reading conservative.

\subsection{The disclosure was ordinary; the attribution was off-scale}
\label{sec:attribution}

\begin{center}
\small
\begin{tabular}{llrl}
\toprule
Date & Event & \texttt{Hugging\_Face} & vs baseline \\\midrule
16 Jul & Hugging Face discloses the breach & 856 & $1.22\times$ \\
17--20 Jul & --- & 439--1{,}144 & flat \\
21 Jul & OpenAI acknowledges responsibility & 1{,}767 & $2.5\times$ \\
22 Jul & press wave & \textbf{17{,}306} & $\mathbf{24.7\times}$ \\
\bottomrule
\end{tabular}

\smallskip
\noindent\footnotesize Baseline 701.5 views/day throughout, per §3.2.
\end{center}

Against the distribution of daily ratios for this page over the preceding 166
days (median $1.00$, p90 $1.21$, maximum $1.55$), \textbf{$1.22\times$ is the
91st percentile} --- 15 of 166 days reach it. The disclosure produced a
detectable, top-decile, entirely ordinary day. The contrast is therefore between
two kinds of ordinary: the victim's own disclosure of a serious breach moved its
page to a level it reaches 9\,\% of days anyway, while attribution to a named
frontier lab moved it to $24.7\times$ baseline, roughly sixteen times beyond the
largest daily ratio in the preceding 166 days. The developer's page stayed at
0.35 baseline-days (§\ref{sec:counting}).

\subsection{At least three attention impulses, and the largest is unrelated}
\label{sec:impulses}

\begin{center}
\small
\begin{tabular}{llrl}
\toprule
Date & Driver & Views & vs baseline \\\midrule
22 Jul & attribution $+$ press wave & 17{,}306 & $24.7\times$ \\
27 Aug & two forensic reports \textbf{and} acquisition reporting & 28{,}026 & $39.9\times$ \\
\textbf{3 Sep} & \textbf{Nvidia--Hugging Face deal announced} & \textbf{34{,}679} & $\mathbf{49.4\times}$ \\
\bottomrule
\end{tabular}

\smallskip
\noindent\footnotesize All ratios against the same baseline of 701.5 views/day
(30-day median ending seven days before 21 July), per §3.2.
\end{center}

The largest spike in the series is an acquisition announcement, not the
incident, and it is genuine human traffic: \texttt{agent=user} 34{,}679 against
\texttt{all-agents} 35{,}898, split desktop 17{,}080 / mobile-web 16{,}673 /
mobile-app 926. The 27 August impulse is not attributable to the forensic
reports either. \emph{The Information} reported that Nvidia had agreed to buy
Hugging Face for \$12.9\,bn on the evening of 26 August, and TechCrunch carried
it at 06:32 UTC on \textbf{27 August} --- the spike day, in the UTC days
Wikimedia reports. The largest corporate story in the company's history lands on
the same pageview day as the two reports, and this series cannot separate them,
so \textbf{only the 22 July impulse is cleanly attributable to the incident} and
whether forensic publication drives a second peak is not testable here.

This is the sharpest available statement of the entity channel's construct
problem. \texttt{Hugging\_Face} measures attention to a \emph{company}, and over a
long enough horizon company news dominates incident news, so any decay analysis
on such a page must be truncated before the next corporate event --- which the
analyst cannot know in advance. A single-impulse collective-memory model
therefore cannot represent this series; the right family is a self-exciting
process with a time-varying exogenous input
\cite{crane2008robust,rizoiu2016hip}, and none is fitted here
(§\ref{sec:limits}).

\subsection{A cross-channel decay comparison that does not survive its null}
\label{sec:failed}

Fitting the three channels separately gives half-lives of $7.05$ hours (Hacker
News comments), $5.45$ days (Wikipedia) and $6.11$ days (GDELT) --- a
$20.8\times$ spread reading as evidence that attention runs on separate clocks
and that the field read the fast one. Two tests remove that reading.

\paragraph{The engagement half-life is a property of the platform.} Thirty
comparable front-page threads, fitted identically:

\begin{center}
\small
\begin{tabular}{lrrrrr}
\toprule
& min & p25 & \textbf{median} & p75 & max \\\midrule
Null half-life (hours), $n=30$ threads & 5.76 & 6.88 & \textbf{7.69} & 8.80 & 16.94 \\
\bottomrule
\end{tabular}

\smallskip
\textbf{Incident thread: $7.05$ h $\rightarrow$ 33rd percentile}
(rank 10 of 30; Clopper--Pearson 95\,\% interval 17.3rd--52.8th).
\end{center}

The incident's thread decays \emph{more slowly} than the median comparable
thread. Every front-page thread in the sample turns over in 6--9 hours, exactly
as \cite{wu2007novelty} reported in 2007. The number measures Hacker News. Part
of the cross-channel spread is also binning: the engagement channel was fitted
hourly and the others daily, and a daily series cannot resolve a sub-day
half-life. Re-binning the same comments to daily gives $0.759$ days rather than
$0.294$, cutting the ratio against the lookups channel from $18.6\times$ to
$7.2\times$ --- a bound on the artifact rather than a measurement of it, since
the daily fit rests on 8 heavily front-loaded bins.

\paragraph{What remains.} Two statements survive. The channels are \emph{not}
shown to differ: a bootstrap on the ratio of the lookups and media half-lives
gives $1.120$, 95\,\% interval $(0.788,\ 1.740)$, which contains 1. And,
independently of this event, \emph{every} front-page aggregator thread turns
over within about eight hours --- so a conversion scorecard read at 48 hours is
always reading a community channel that has already closed. That is a statement
about the timing of measurement, weaker and more general than a claim about this
warning shot's decay.

\subsection{Robustness}
\label{sec:robust}

Every surviving number was re-tested against the strongest objections we could
construct: log-space OLS bias refit by raw-scale nonlinear least squares,
autocorrelation absorbed by a moving-block bootstrap, contamination probed by
dropping the Kimi-K3 release days, a searched breakpoint charged $2\ln C$ in
AIC, the concept null rebuilt season-matched, the reach comparison recomputed in
both unit systems, half-life comparisons bootstrapped as ratios rather than read
off overlapping intervals, and the engagement half-life set against 30
comparable threads. Magnitudes move --- up to $2.7\times$ under a change of
estimator --- but every pairwise verdict reported here survives, and the two
claims that do not survive are the ones withdrawn in §\ref{sec:failed}.
App.~\ref{sec:robust-table} gives the objection-by-objection table, and
\texttt{scripts/robustness/} reproduces it in one command.

%====================================================================
\section{Discussion}
\label{sec:discussion}
%====================================================================

\paragraph{The finding a communications team can act on.} Before asking how far
an incident travelled, decide whose page counts as ``it'', and say so. On an
incident with a perpetrator and a victim the available choices give opposite
answers, and for short windows some candidates do not exist yet --- the
difference between ``the warning shot underperformed an export-control action''
and ``the warning shot outperformed it'', from the same data.

\paragraph{The result that should worry the field.} The organisations moved; the
vocabulary shows no movement we could detect. Concept-page attention over 30
days is indistinguishable from what those pages do in an ordinary summer, and
the generic control sits inside the same band. The honest form is bounded --- the
design could not have detected a response smaller than about $2.2\times$ ordinary
drift, so what can be said is that no \emph{large} conversion occurred. Even
bounded that way it is the result most worth acting on, because a large
conversion is what the warning-shot argument needs.

\paragraph{On measurement timing.} Every front-page aggregator thread closes
within about eight hours \cite{wu2007novelty}, so a conversion scorecard read at
48 hours is reading a channel that shut a day and a half earlier --- a property
of the medium rather than of any event, and what the cross-channel analysis
supports once its null is applied (§\ref{sec:failed}).

\paragraph{Two hypotheses, not rules.} From one incident: (i) plan for the
attention event to fire on attribution to a named lab rather than on the
victim's disclosure; (ii) expect corporate news about the victim to overwhelm
incident signal on any long horizon. Each is generated by $n=1$. Both concern
\emph{whom} and \emph{when} rather than volume --- as does the strongest
communications finding from cabinet-level AI crisis exercises, that ``specific,
practical information for the most-exposed actors was more useful than broad
warnings or silence'' \cite{shamir2026ttx}.

%====================================================================
\section{Limitations and Dual-Use Considerations}
\label{sec:limits}
%====================================================================

\paragraph{Limitations.} (L1) \emph{$n = 1$ event}, with one matched control:
nothing here establishes a law of warning-shot attention, and
§\ref{sec:discussion}'s implications are hypotheses from a single incident.
(L2) \emph{Proxy validity.} Pageviews measure encyclopaedia lookups, not
awareness or opinion change; \texttt{Hugging\_Face} measures a \emph{company},
and §\ref{sec:impulses} shows corporate news dominating it; GDELT
\texttt{timelinevol} is a share of monitored coverage, not a count.
(L3) \emph{Small nulls.} Thirteen overlapping concept windows collapse to a
handful of independent samples, so the 53.8th percentile is a location inside
ordinary variation, not a $p$-value; the thread null is 30 threads from one
platform over three months. (L4) \emph{Estimator dependence.} Half-lives move up
to $2.7\times$ between log-space OLS and raw-scale NLS, so no single half-life is
the quantity. (L5) \emph{Truncation is unavoidable and partly arbitrary}: decay
windows must end before the next corporate event, which is not knowable in
advance. (L6) \emph{Confounded second impulse}: the 27 August spike has three
candidate drivers, none separable here, so whether forensic publication drives a
second peak is untestable on these data. (L7) \emph{Comparator dependence of the
headline}: the inversion's direction is robust across all six June candidates
and both unit systems, but its magnitude spans $3.9\times$ to $145.5\times$
normalised and $1.46\times$ to $92.5\times$ raw, and no minimum-traffic rule was
pre-registered. (L8) \emph{The selection rule is retrospective}: ``count the
absorbing page'' is an argmax over the outcome, a diagnostic rather than an
estimator. (L9) \emph{Power}: the concept null could not have detected a response
below about $2.2\times$ ordinary drift. (L10) \emph{No self-exciting model
fitted}, though \cite{crane2008robust,rizoiu2016hip} is the right family here.
(L11) \emph{Review provenance}: the adversarial review behind this paper was the
author's own with model assistance, not an independent party's, so it tests the
objections the author could construct. App.~\ref{sec:limits-ext} expands L3, L4,
L6, L7 and L10.

\paragraph{Dual-use.} This work measures what makes an AI incident travel, and
the same findings could be used to suppress attention rather than spread it.
§\ref{sec:attribution} is the sharp case: a victim's unilateral disclosure moves
its page to an ordinary busy day, while attribution to a named lab moves it
off-scale. A lab wishing to minimise public attention to its own incident could
read that as an argument for delaying or diffusing attribution, and
§\ref{sec:counting} as an argument for steering coverage toward whichever entity
page is least likely to register it.

We state this openly for three reasons. It is already the incentive gradient labs
face, so the measurement describes a pressure that exists rather than creating
one. A public, checkable estimate is more useful to those arguing for timely
attribution --- regulators, journalists, the breached party --- than to the one
actor deciding whether to delay, because the number that shows delay is
effective also shows when it has happened. And it cuts symmetrically, being
equally an argument for a regime in which attribution is timely by requirement
rather than by choice, which is the use we would defend. No part of this
analysis uses non-public data, touches any live system, or involves any model
capability; the released artifact holds public pageview counts, thread
timestamps and coverage volumes, and nothing in it lowers the cost of running an
incident or identifies a person.

\bibliographystyle{plainurl}
\bibliography{references}

%====================================================================
\appendix
%====================================================================

\section{Extended limitations}
\label{sec:limits-ext}

\paragraph{(L3) Null sizes.} The concept null is thirteen 30-day windows slid in
5-day steps across 15 June--15 September of 2024 and 2025. Overlapping windows
are not independent, so the effective sample is smaller than thirteen. The
observed value is rank 7 of 13; an exact binomial (Clopper--Pearson) interval on
that rank runs from the \textbf{25.1st to the 80.8th percentile}, and the true
interval is wider still because of the overlap. The percentile is therefore a
location statement inside ordinary variation, not a $p$-value, and ``53.8''
should not be read as a three-significant-figure quantity. The thread null is 30
threads drawn from one platform over three months.

\paragraph{(L4) Estimator dependence.} Least squares on log-transformed data is
a biased estimator with invalid standard errors \cite{clauset2009powerlaw}.
Every exponential fit is therefore also computed by nonlinear least squares on
the raw scale. Half-lives on the daily channels move by up to $2.7\times$
between the two, and the cross-channel spread moves from $20.8\times$ to
$12.2\times$; the ordering of the channels is invariant. No single half-life in
this paper should be treated as the quantity.

\paragraph{(L6) The confounded second impulse.} The 27 August spike has three
candidate drivers: OpenAI's technical report, the independent METR/Redwood
report \cite{metr2026investigation}, and reporting of the Nvidia acquisition.
Business Insider reported takeover interest over the preceding weekend;
\emph{The Information} reported the \$12.9\,bn agreement on the evening of
26 August; TechCrunch carried it at 23:32 PDT that day, which is 06:32 UTC on
27 August --- the spike day, in the UTC days Wikimedia reports. None of the
three is separable in this series.

\paragraph{(L7) Comparator dependence, decomposed.} A ratio of baseline-days
factors exactly into a raw-excess ratio times an inverse-baseline ratio,
$(e_1/b_1)/(e_2/b_2) = (e_1/e_2)\times(b_2/b_1)$:

\begin{center}
\small
\begin{tabular}{lrrr}
\toprule
Comparator & Baseline-days ratio & $=$ raw ratio & $\times$ baseline ratio \\\midrule
\texttt{Anthropic} (12 Jun anchor) & $28.92\times$ & $1.78\times$ & $16.28\times$ \\
\texttt{Anthropic} (13 Jun anchor) & $23.06\times$ & $1.46\times$ & $15.82\times$ \\
\texttt{Claude\_(language\_model)}  & $91.20\times$ & $8.15\times$ & $11.19\times$ \\
\texttt{OpenAI}                    & $145.53\times$ & $13.40\times$ & $10.86\times$ \\
\texttt{Export\_control}           & $3.89\times$ & $92.51\times$ & $0.04\times$ \\
\bottomrule
\end{tabular}
\end{center}

\noindent Most of the \texttt{Anthropic} multiple is \texttt{Anthropic} being a
sixteen-times larger page, not the July event drawing sixteen times more
traffic. The \emph{smallest} multiple in the range, $3.89\times$ against
\texttt{Export\_control}, conceals a raw-excess ratio of $92.51\times$ running
the other way, because \texttt{Export\_control}'s baseline is a twenty-fifth of
the victim page's. Neither the largest nor the smallest multiple should be
quoted as the result.

\paragraph{(L10) No self-exciting model fitted.} A self-exciting process with a
time-varying exogenous input \cite{crane2008robust,rizoiu2016hip} is the right
family for a series with multiple impulses. A reference HIP implementation was
run and did not identify on a 50-day series with two sparse impulses; a denser
exogenous input is the next step. The repository ships the HawkesN fit and the
feasibility probe so the negative result is checkable.

\paragraph{Citation discipline.} Every reference had its title, authors and year
confirmed against the source. A German-language headline widely quoted in
discussion of this incident could not be located and is therefore not cited.

\section{Robustness table}
\label{sec:robust-table}

Each row is an objection to a number in §\ref{sec:counting}--§\ref{sec:failed},
the test run against it, and the outcome. \texttt{scripts/robustness/} reproduces
every row offline.

\begin{center}
\small
\begin{tabular}{p{0.29\linewidth}p{0.30\linewidth}p{0.31\linewidth}}
\toprule
Objection & Test & Outcome \\\midrule
Log-space OLS is biased \cite{clauset2009powerlaw} & refit by raw-scale NLS & magnitudes move up to $2.7\times$; ordering invariant \\
Residuals autocorrelated & moving-block bootstrap & intervals widen 3--18\,\%; all pairwise verdicts unchanged \\
Entity page contaminated & drop the Kimi-K3 release days & OLS half-life $5.450 \to 5.413$ d ($0.99\times$) \\
Breakpoint chosen by search & charge $2\ln C$ to two-phase AIC & two of three channels survive, not three \\
Null not season-matched & rebuild from prior-year Jul--Sep & $68.2 \to \mathbf{53.8}$; matched figure reported \\
Comparison uses raw views & recompute in baseline-days & $1.46\times$ raw, $\mathbf{23.06\times}$ normalised; both reported \\
Overlapping intervals used as a test & bootstrap the ratio directly & $(0.788,\ 1.740)$, contains 1 \\
Half-life may be platform-typical & null over 30 comparable threads & \textbf{33rd percentile --- claim not supported} \\
\bottomrule
\end{tabular}
\end{center}

\section{Search protocol for the novelty claim}
\label{sec:searchprotocol}

The claim in §\ref{sec:related} is narrow: no prior treatment of
\emph{entity-choice dependence} in incident reach measurement could be found.
Searched 2026-09-04/05 via Google Scholar and web search, in English, with
combinations of \{collective attention, attention decay, Wikipedia pageviews\}
$\times$ \{AI incident, data breach, cybersecurity incident, security
disclosure\}, plus forward citations of \cite{garciagavilanes2016crashes} and
\cite{igarashi2022twophase}. This is a bounded search, not a systematic review.
The works in §\ref{sec:related} that most closely bound the contribution,
\cite{cavusoglu2004breach} among them, were located only after external
prompting --- which is itself evidence that the search was not exhaustive.

\section{Reproducibility}
\label{sec:repro}

Python 3.10+; \texttt{numpy}, \texttt{scipy}, \texttt{matplotlib},
\texttt{requests}, \texttt{pyyaml}. Bootstrap seeds are fixed (block bootstrap
seed 1, ratio bootstrap seed 4).

\begin{center}
\small
\begin{tabular}{ll}
\toprule
Command & What it establishes \\\midrule
\texttt{python -m pytest} & 115 tests, no network \\
\texttt{python scripts/verify.py} & 121 checks recomputed from committed results \\
\texttt{python scripts/run.py all} & regenerates the channel fits and figures \\
\texttt{python scripts/robustness\_suite.py} & regenerates the robustness table \\
\texttt{python scripts/nulls\_and\_controls.py} & regenerates the thread and season-matched nulls \\
\texttt{python scripts/power\_and\_comparators.py} & regenerates the power bound and comparator sweep \\
\bottomrule
\end{tabular}
\end{center}

\noindent Every command runs offline against the shipped cache at \$0 cost. A
clean checkout without the sibling corpora skips the corpus tests.

\section{Declarations}

\paragraph{Data availability.} All sources are public and keyless: Wikimedia
Pageviews \cite{wikimedia2026pageviews}, the Hacker News Algolia API, and GDELT
DOC 2.0 \cite{gdelt2026doc}. Code, cached API responses, committed results and
figures accompany this submission and are released MIT-licensed. No restricted,
proprietary or personal data was used.

\paragraph{Ethics.} No human subjects, no personal data, no interaction with any
live system beyond public read-only API queries issued with an identifying user
agent and rate-limit backoff. No journalist, creator or policymaker was
contacted: the sprint brief conditions outreach on affiliation with a
credibility-lending structure, and that condition is not met here.

\paragraph{Author contributions (CRediT).} Fatimah Emad Eldin:
conceptualization, methodology, software, formal analysis, investigation, data
curation, writing---original draft, writing---review \& editing, visualization.

\paragraph{Conflict of interest.} None declared. The author has no affiliation
with OpenAI, Hugging Face, Anthropic, Nvidia, METR, Redwood Research, CeSIA,
RAND or Apart Research. §\ref{sec:counting} re-measures figures published by
CeSIA, a co-organiser of this sprint; the re-measurement is reported as a
methodological result with the benign explanation stated alongside it.

\paragraph{Funding.} None. No compute budget and no paid API access were used.

\paragraph{LLM-assisted authoring.} The author used Claude Opus for literature
scoping, API feasibility probing, prose drafting, code review and adversarial
review of the drafts. The research question, the entity-choice framing and the
null designs are the author's own. All quantitative results were computed by the
released code and are reproducible via \texttt{verify.py}.

\end{document}
